\documentclass[11pt,a4paper]{article}
\usepackage[italian]{babel}
\usepackage[margin=2.5cm]{geometry}
\usepackage{parskip}
\usepackage{amsmath, amssymb, amsthm}
\usepackage[most]{tcolorbox}
\usepackage{hyperref}

% --- Riquadri con bordo nero, senza numero ---
\tcbset{nota/.style={colback=white, colframe=black, boxrule=0.5pt,
  sharp corners}}
\NewTColorBox{definizione}{}{nota,
  before upper={\textbf{Definizione.}\ }}
\NewTColorBox{teorema}{o}{nota,
  before upper={\textbf{Teorema\IfValueT{#1}{ (#1)}.}\ }}

% --- Ambienti semplici (senza riquadro) ---
\theoremstyle{definition}
\newtheorem*{esempio}{Esempio}
\newtheorem*{osservazione}{Osservazione}

\title{Principio di induzione}
\author{Marco Cerbella}
\date{Lezione 4 -- 29 settembre 2026}

\begin{document}
\maketitle

\section{Il principio di induzione}

\begin{osservazione}
Sia $U \subseteq \mathbb{N}$ tale che
\begin{enumerate}
    \item $0 \in U$;
    \item $k \in U \Rightarrow k + 1 \in U$.
\end{enumerate}
Allora $U = \mathbb{N}$.
\end{osservazione}

Infatti $0 \in U$, quindi $1 \in U$, quindi $2 \in U$, e così via: ``si sale
a gradini'' e si raggiungono tutti i naturali. Se $p(n)$ è un predicato
sui naturali, si applica questa idea all'insieme $U$ dei $n$ per cui
$p(n)$ è vero.

\begin{teorema}[Principio di induzione]
Sia $p(n)$ un predicato. Supponiamo che:
\begin{enumerate}
    \item $p(0)$ sia vero \hfill [\emph{base dell'induzione}];
    \item $\forall\, k \ \big(p(k) \text{ vero} \Rightarrow p(k+1)
          \text{ vero}\big)$ \hfill [\emph{passo induttivo}].
\end{enumerate}
Allora $\forall\, n \in \mathbb{N}$ $p(n)$ è vero.
\end{teorema}

Nel passo induttivo l'ipotesi $p(k)$ si chiama \emph{ipotesi induttiva} e
$p(k+1)$ è la \emph{tesi da dimostrare}.

\begin{esempio}
Dimostrare che per ogni $n$ naturale $2 \mid n(n+1)$.

Sia $p(n)$ il predicato ``$2 \mid n(n+1)$''.

\textbf{1) Caso base.} $p(0)$: $2 \mid 0 \cdot (0 + 1) = 0$. Infatti
$2 \mid 0$ significa $0 = k \cdot 2$ con $k \in \mathbb{Z}$, ed è vero per
$k = 0$: $0 = 0 \cdot 2$.

\textbf{2) Passo induttivo.} Dobbiamo dimostrare $p(k) \Rightarrow
p(k+1)$.
\begin{itemize}
    \item \emph{Ipotesi induttiva} $p(k)$: $2 \mid k(k+1)$.
    \item \emph{Tesi da dimostrare} $p(k+1)$: $2 \mid (k+1)\big((k+1)+1\big)$.
\end{itemize}
Calcoliamo
\[
(k+1)(k+2) = (k+1)k + 2(k+1).
\]
Per l'ipotesi induttiva $2 \mid k(k+1)$, cioè $k(k+1) = 2m$ con
$m \in \mathbb{Z}$. Quindi
\[
(k+1)(k+2) = 2m + 2(k+1) = 2\,(m + k + 1),
\]
cioè $(k+1)(k+2)$ è multiplo di $2$, ovvero $2 \mid (k+1)(k+2)$, che è la
tesi da dimostrare.
\end{esempio}

\section{Induzione a partire da $\bar n$}

Se il predicato è vero solo da un certo naturale in poi, si parte da
quel punto.

\begin{teorema}
Sia $p(n)$ una proposizione e sia $\bar n \in \mathbb{N}$. Supponiamo che:
\begin{enumerate}
    \item $p(\bar n)$ sia vero;
    \item $\forall\, k \geq \bar n \ \big(p(k) \text{ vero} \Rightarrow
          p(k+1) \text{ vero}\big)$.
\end{enumerate}
Allora $\forall\, n \geq \bar n$ $p(n)$ è vera.
\end{teorema}

\begin{esempio}
Il predicato $p(n)$: ``$n - 5 > 0$'' è vero da $6$ in poi: si applica il
teorema con $\bar n = 6$.
\end{esempio}

\begin{esempio}
Dimostrare che per ogni $n \geq 1$
\[
1 + 2 + \dots + n = \frac{n(n+1)}{2}. \tag*{$p(n)$}
\]

\textbf{1) Caso base} $\bar n = 1$. $p(1)$: $1 = \dfrac{1 \cdot (1+1)}{2}$,
vera.

\textbf{2) Passo induttivo.}
\begin{itemize}
    \item \emph{Ipotesi induttiva} $p(k)$:
          $1 + \dots + k = \dfrac{k(k+1)}{2}$.
    \item \emph{Tesi da dimostrare} $p(k+1)$:
          $1 + \dots + k + (k+1) = \dfrac{(k+1)(k+2)}{2}$.
\end{itemize}
Si ha
\begin{align*}
1 + 2 + \dots + k + (k+1) &= (1 + \dots + k) + (k+1)\\
&= \frac{k(k+1)}{2} + (k+1) \qquad \text{(per ipotesi induttiva)}\\
&= \frac{k(k+1) + 2(k+1)}{2} = \frac{(k+1)(k+2)}{2},
\end{align*}
che è la tesi. Per induzione $p(n)$ vale per ogni $n \geq 1$.
\end{esempio}

\end{document}
